\ncertchapter{1}{Measurement, Units and Uncertainty}\label{ch:measurement}

Every number in physics is a measurement, and every measurement carries two things that a bare number
hides: a \emph{unit}, which says what was counted, and an \emph{uncertainty}, which says how well. A
student who loses track of either will eventually produce an answer that is wrong by a factor of a
thousand, or one that is quoted to six figures when the data justified two. Both errors are common,
both are avoidable, and both are habits rather than pieces of knowledge.

This chapter builds three such habits. The first is dimensional bookkeeping: writing every quantity in
terms of mass, length and time, so that a formula can be checked before it is trusted and, sometimes,
derived before it is understood. The second is estimation, the ability to know roughly what an answer
should be before calculating it, which is the cheapest error check there is. The third is honest
uncertainty: knowing how errors in the inputs combine into an error in the result, and therefore which
measurement is worth improving.

None of this is glamorous. It is placed first in the physics because every
later calculation leans on it, and because the examinations reward it reliably: dimensional analysis
and error propagation are among the most predictable sources of marks in the entire syllabus.

\section{Units and Dimensions}

Physics is quantitative, and every quantity is a \emph{number} times a \emph{unit}: ``\SI{5}{m}''
means five times an agreed standard length. The SI fixes seven \textbf{base units} (kilogram, metre,
second, ampere, kelvin, mole and candela), chosen so that each is independent of the others; every
other unit is a \textbf{derived unit} built from them, as the newton is $\si{kg.m/s^2}$. Three base units
suffice for all of mechanics: mass (kilogram, kg), length (metre, m), and time (second, s).

\begin{defbox}{Dimensions}
The \textbf{dimension} of a quantity records how it is built from mass $\mathsf{M}$, length
$\mathsf{L}$, and time $\mathsf{T}$. We write $[Q]$ for the dimension of $Q$. For example,
$[\text{velocity}]=\mathsf{L}\mathsf{T}^{-1}$, $[\text{force}]=\mathsf{M}\mathsf{L}\mathsf{T}^{-2}$,
$[\text{energy}]=\mathsf{M}\mathsf{L}^2\mathsf{T}^{-2}$.
\end{defbox}

\subsection{Dimensional analysis: the free consistency check}

Dimensions must balance across any valid equation; you cannot add a length to a time any more than
you can add apples to hours. This gives two tools at no cost.

\begin{methodbox}{The two uses of dimensions}
\textbf{To check a formula:} confirm every term has the same dimension. A single mismatch \emph{kills}
the formula. \\[3pt]
\textbf{To find a dependence:} assume $Q=k\,x^a y^b z^c$, match each of $\mathsf{M},\mathsf{L},\mathsf{T}$
separately, and solve for $a,b,c$. You recover the \emph{shape} of the law; never the dimensionless
constant $k$.
\end{methodbox}

\begin{trapbox}{function arguments must be dimensionless; constants usually are not}
Arguments of $\sin$, $\exp$, and $\log$ must be dimensionless. In $\sin(ct)$, with $t$ a time, this forces
$[c]=\mathsf{T}^{-1}$, often the fastest way to find a constant's dimensions. And constants of nature, such as
the gravitational constant $G$, are \emph{not} dimensionless: read their dimensions off an equation you
trust, as is done for $G$ and others in \S\ref{sec:dimformulae}.
\end{trapbox}

\begin{exbox}{The period of a pendulum, without solving anything}
Guess that a pendulum's period $T$ (the time for one full swing, so its dimension is $\mathsf{T}$)
depends on its length $\ell$, its mass $m$, and the acceleration due to gravity $g$ (in \si{m/s^2}, so
$[g]=\mathsf{LT^{-2}}$). Write $T=k\,\ell^a m^b g^c$, with $k$ a pure number, and match dimensions:
\[
\mathsf{T}=\mathsf{L}^a\,\mathsf{M}^b\,(\mathsf{L}\mathsf{T}^{-2})^c
=\mathsf{M}^b\,\mathsf{L}^{a+c}\,\mathsf{T}^{-2c}.
\]
Matching: $\mathsf{M}: b=0$; $\mathsf{T}: -2c=1\Rightarrow c=-\tfrac12$; $\mathsf{L}: a+c=0\Rightarrow
a=\tfrac12$. So $T=k\sqrt{\ell/g}$, the mass drops out, and the length enters as a square root, both
famous facts, obtained without touching a differential equation. \ans{$T\propto\sqrt{\ell/g}$}

\smallskip
Notice three things. We got a testable prediction for free, $T\propto\sqrt\ell$. The mass drops out
because nothing else on the list contains $\mathsf{M}$ for it to pair with; but the conclusion is only as
good as the list, since the amplitude $\theta_0$ (an angle, so dimensionless) could enter as any function
$f(\theta_0)$ unseen, which is why experiment is still needed. And $k$ is out of reach: its true value,
$2\pi$, comes from the full theory.
\end{exbox}

\subsection{The dimensional formulae worth deriving once}\label{sec:dimformulae}

Exam questions ask for the dimensions of specific constants over and over. Do not memorise the list;
memorise the \emph{method}: write a defining equation, replace each symbol by its dimensions, simplify.
Below, that method is applied to the ten quantities that recur most. Most belong to later chapters or
to other branches of physics; all you need now is the defining equation quoted for each, not its
meaning.

\begin{derivbox}{Ten dimensional formulae, from defining equations}
\stepn{1} \textbf{Force.} Newton's second law, force $=$ mass $\times$ acceleration: $F=ma \Rightarrow [F]=\mathsf{M}\cdot\mathsf{LT^{-2}}
=\boxed{\mathsf{M\,L\,T^{-2}}}$ \;(newton).

\stepn{2} \textbf{Work / energy.} $W=Fd\cos\theta$ (force times distance $d$ moved, with $\theta$ the angle between them), and $\cos\theta$ is a pure number, so
$[W]=\boxed{\mathsf{M\,L^{2}\,T^{-2}}}$ \;(joule).

\stepn{3} \textbf{Torque} (turning effect of a force $F$ applied at distance $r$ from a pivot). $\tau=rF\sin\theta \Rightarrow [\tau]=\boxed{\mathsf{M\,L^{2}\,T^{-2}}}$.
\emph{Identical to work}, yet torque has a direction (it is a vector, \chref{ch:vectors}) and work does not. Dimensions cannot always tell
quantities apart.

\stepn{4} \textbf{Gravitational constant.} From $F=\dfrac{Gm_1m_2}{r^{2}}$, $G=\dfrac{Fr^{2}}{m_1m_2}$:
\[
[G]=\frac{\mathsf{M\,L\,T^{-2}}\cdot\mathsf{L^{2}}}{\mathsf{M^{2}}}
=\boxed{\mathsf{M^{-1}\,L^{3}\,T^{-2}}}.
\]

\stepn{5} \textbf{Coefficient of viscosity.} Newton's law of viscous flow is $F=\eta A\dfrac{dv}{dy}$, and
the velocity gradient has dimensions $\mathsf{LT^{-1}/L}=\mathsf{T^{-1}}$, so
\[
[\eta]=\frac{\mathsf{M\,L\,T^{-2}}}{\mathsf{L^{2}}\cdot\mathsf{T^{-1}}}
=\boxed{\mathsf{M\,L^{-1}\,T^{-1}}}\quad(\si{Pa.s}).
\]

\stepn{6} \textbf{Moment of inertia} (a sum over the masses $m$ of a body, each at distance $r$ from an axis). $I=\sum mr^{2}\Rightarrow[I]=\boxed{\mathsf{M\,L^{2}}}$, no time
appears at all.

\stepn{7} \textbf{Planck constant.} $E=h\nu$, the energy $E$ of a quantum of light of frequency $\nu$
(oscillations per second, so $[\nu]=\mathsf{T^{-1}}$):
\[
[h]=\frac{\mathsf{M\,L^{2}\,T^{-2}}}{\mathsf{T^{-1}}}=\boxed{\mathsf{M\,L^{2}\,T^{-1}}}.
\]
These are the dimensions of \emph{action} (energy $\times$ time), which is why $h$ is called the quantum
of action. They are also the dimensions of angular momentum (\chref{ch:rotation}), and that is no
coincidence: in quantum physics angular momentum comes in units of $h/2\pi$.

\stepn{8} \textbf{Universal gas constant.} Temperature and amount of substance are base quantities in
their own right, and SI writes their dimensions as $\Theta$ and $\mathsf{N}$. From the ideal-gas law
$PV=nRT$ (pressure $P$, volume $V$, amount $n$ in moles, absolute temperature $T$), and noting
$PV$ is an energy,
\[
[R]=\frac{\mathsf{M\,L^{2}\,T^{-2}}}{\mathsf{N}\,\Theta}
=\boxed{\mathsf{M\,L^{2}\,T^{-2}}\,\Theta^{-1}\,\mathsf{N}^{-1}}.
\]

\stepn{9} \textbf{Boltzmann constant.} $k_B=R/N_A$, the gas constant \emph{per molecule}, where the Avogadro constant $N_A$ is
the number of molecules in a mole. Dropping the
mole leaves $[k_B]=\boxed{\mathsf{M\,L^{2}\,T^{-2}}\,\Theta^{-1}}$.

\stepn{10} \textbf{Stefan constant.} Stefan--Boltzmann law $\dfrac{P}{A}=\sigma T^{4}$ (power $P$ radiated per area $A$ by a
black body at absolute temperature $T$), with power
$[P]=\mathsf{M L^{2}T^{-3}}$:
\[
[\sigma]=\frac{\mathsf{M\,L^{2}\,T^{-3}}}{\mathsf{L^{2}}\cdot\Theta^{4}}
=\boxed{\mathsf{M\,T^{-3}}\,\Theta^{-4}}.
\]
\end{derivbox}

\begin{tablebox}{Reference table: check your derivation against this, do not memorise it}
\begin{center}
\begin{tabular}{@{}llll@{}}
\toprule
Quantity & Defining relation & Dimensional formula & SI unit\\
\midrule
Force $F$              & $ma$                  & $\mathsf{M\,L\,T^{-2}}$                          & \si{N}\\
Work / energy $W$      & $Fd\cos\theta$        & $\mathsf{M\,L^{2}\,T^{-2}}$                      & \si{J}\\
Torque $\tau$          & $rF\sin\theta$        & $\mathsf{M\,L^{2}\,T^{-2}}$                      & \si{N.m}\\
Grav.\ constant $G$    & $Fr^{2}/m_1m_2$       & $\mathsf{M^{-1}\,L^{3}\,T^{-2}}$                 & \si{N.m^2/kg^2}\\
Viscosity $\eta$       & $F/[A(dv/dy)]$        & $\mathsf{M\,L^{-1}\,T^{-1}}$                     & \si{Pa.s}\\
Moment of inertia $I$  & $\sum mr^{2}$         & $\mathsf{M\,L^{2}}$                              & \si{kg.m^2}\\
Planck constant $h$    & $E/\nu$               & $\mathsf{M\,L^{2}\,T^{-1}}$                      & \si{J.s}\\
Gas constant $R$       & $PV/nT$               & $\mathsf{M\,L^{2}\,T^{-2}}\,\Theta^{-1}\mathsf{N}^{-1}$ & \si{J/(mol.K)}\\
Boltzmann $k_B$        & $R/N_A$               & $\mathsf{M\,L^{2}\,T^{-2}}\,\Theta^{-1}$          & \si{J/K}\\
Stefan $\sigma$        & $P/(AT^{4})$          & $\mathsf{M\,T^{-3}}\,\Theta^{-4}$                 & \si{W/(m^2.K^4)}\\
\bottomrule
\end{tabular}
\end{center}
\end{tablebox}

\begin{defbox}{The principle of homogeneity of dimensions}
The rule underlying everything above deserves its name:
\begin{center}
\emph{Every term that is added, subtracted or equated in a physical equation must have the same
dimensions.}
\end{center}
It is really a statement that our choice of units is arbitrary: an equation must stay true when we
switch from metres to feet, and only a dimensionally homogeneous equation survives that switch. Its four
uses are checking an equation, deriving the \emph{form} of a law, converting between unit systems, and
deducing an unknown constant's dimensions.
\end{defbox}

\begin{methodbox}{Converting units with dimensions}
A physical quantity does not change when the unit changes; only the number does. If the quantity has
dimensions $\mathsf{M}^a\mathsf{L}^b\mathsf{T}^c$, its value $n_1$ in one system and $n_2$ in another
satisfy $n_1u_1=n_2u_2$, where $u_1$ and $u_2$ are the two units, so
\[
n_2=n_1\Big(\frac{M_1}{M_2}\Big)^{a}\Big(\frac{L_1}{L_2}\Big)^{b}\Big(\frac{T_1}{T_2}\Big)^{c},
\]
where $M_1/M_2$ is how many of the new mass units make one old one, and so on. A bigger unit gives a
smaller number. For example, energy ($a=1$, $b=2$, $c=-2$) goes from SI to CGS with
$M_1/M_2=\SI{1}{kg}/\SI{1}{g}=10^{3}$ and $L_1/L_2=\SI{1}{m}/\SI{1}{cm}=10^{2}$, so
$\SI{1}{J}=1\times10^{3}\times(10^{2})^{2}=10^{7}$ erg. The everyday case is the same rule:
$\SI{72}{km/h}=72\times\frac{1000}{3600}\,\si{m/s}=\SI{20}{m/s}$.
\end{methodbox}

\subsection{What dimensional analysis cannot do}

The method is a filter, not an oracle, and knowing exactly where its power stops is what makes it usable.
It fails in four specific ways:

\begin{trapbox}{the four blind spots}
\textbf{1. It cannot find dimensionless factors.} $T=2\pi\sqrt{\ell/g}$ and $T=\sqrt{\ell/g}$ are
dimensionally identical; the $2\pi$ must come from real physics. You get the \emph{shape}, never the
number in front.\\[3pt]
\textbf{2. It cannot handle sums.} In $s=ut+\tfrac12at^2$ both terms have dimension $\mathsf L$; nothing
in the method tells you both are needed, or that the coefficient is $\tfrac12$.\\[3pt]
\textbf{3. It is blind to dimensionless variables.} An angle, a mass ratio, or the Reynolds number can
appear raised to \emph{any} power without disturbing the dimensions, so the method cannot see it.\\[3pt]
\textbf{4. It cannot distinguish quantities with the same dimensions.} Work and torque are both
$\mathsf M\mathsf L^2\mathsf T^{-2}$, yet one is a scalar and one a vector, and adding them is nonsense.
\end{trapbox}

But what it \emph{can} do, it does absolutely: a dimensionally wrong equation is not approximately wrong,
not wrong in some regime; it is not physics. So check every answer you derive. It costs five seconds
and catches algebra slips that re-reading often misses.

\why{Dimensional analysis is the physicist's first line of defence. In an exam it catches a wrong
formula in seconds; in research it has repeatedly given the form of a law before the detailed theory was worked out. The
best-known case is G.~I.~Taylor's estimate of the energy released by the first atomic explosion, made
from published photographs of the growing fireball while the official figure was still secret.}

\section{Estimation and Significant Figures}

Dimensions tell us whether an answer has the right \emph{form}; they say nothing about whether its
\emph{size} is sensible, or how many of its digits deserve belief. Those are the two questions of this
section.

\subsection{Order-of-magnitude thinking}

Physicists estimate constantly, and it is a trained skill, not a talent. The claim: you can land within a
factor of ten of almost anything by decomposing it into pieces you can guess, because independent errors
tend to cancel: overestimate one factor, underestimate the next, and the product lands close.

\begin{tablebox}{Anchors worth memorising}
\centering
\begin{tabular}{@{}ll@{\hspace{2.2em}}ll@{}}
\toprule
Adult height & \SI{1.7}{m} & Density of water & \SI{1000}{kg/m^3}\\
Earth's radius & \SI{6400}{km} & Atomic size & \SI{e-10}{m}\\
$g$ & $\SI{9.8}{m/s^2}\approx10$ & Avogadro $N_A$ & \SI{6e23}{mol^{-1}}\\
Atmospheric pressure & \SI{e5}{Pa} & Speed of sound (air) & \SI{340}{m/s}\\
\bottomrule
\end{tabular}
\end{tablebox}

The value of an estimate is not the number; it is the \emph{veto}. An estimate tells you instantly when
an answer is absurd: a car of mass \SI{3}{kg}, or an electron at \SI{e12}{m/s}, is a slip you catch
without rechecking the algebra.

\begin{exbox}{How many breaths in a lifetime}
Decompose into guessable pieces: breaths/min $\approx15$, minutes/year $\approx5\times10^5$, lifetime
$\approx75$ years. Then
\[
N\approx 15\times(5\times10^5)\times75\approx5.6\times10^8\approx10^9.
\]
About a billion. How much to trust it? Each input might be off by 50\%, but the three are independent,
and independent errors combine \emph{in quadrature} (as the square root of the sum of squares, a rule
met again in \S\ref{sec:propagation}), not by adding: $\sqrt3\times50\%\approx90\%$,
roughly a factor of two, rather than the worst case $1.5^{3}\approx3.4$ in which all three err the same
way. Good to an order of magnitude, exactly as claimed. \ans{$\sim10^9$}
\end{exbox}

\subsection{Significant figures}

Your calculator will hand you $9.813725490196078$. Writing that down is not precision; it is a lie about
your apparatus. Significant figures are a crude way of tracking how much of your answer is knowledge and
how much is arithmetic noise. A reading written as $2.5$ claims that the last digit is roughly
right and nothing beyond it is known; $2.50$ claims ten times more. The significant figures are the
digits that carry such a claim. Counting them follows from that meaning: every non-zero digit counts,
and so does a zero between them ($2.05$ has three); leading zeros only place the decimal point and never
count ($0.0052$ has two); trailing zeros after a decimal point do count ($2.50$ has three). A trailing
zero in a whole number is ambiguous ($1200$), so write $1.20\times10^{3}$ to claim three. Two rules for
arithmetic, and two habits, cover most remaining cases:

\begin{itemize}
\item In \textbf{multiplication/division}, the result carries as many significant figures as the
\emph{least precise} input. $(2.3)(4.5678)=11$ (or $1.1\times10^{1}$): the least precise input has two significant figures, so
the answer has two; not $10.50594$, and not $10.5$.
\item In \textbf{addition/subtraction}, match the least precise \emph{decimal place}, not the figure
count: $12.11+0.3=12.4$.
\item Keep \textbf{guard digits} through intermediate steps and round only at the end, or rounding errors
compound.
\item Exact numbers (a factor of 2, a count of 5 objects) have infinite significant figures and
never limit anything.
\end{itemize}

\section{Uncertainty and Its Propagation}

Significant figures only hint at how well a number is known. To say it properly we attach an explicit
uncertainty to each measurement and work out how it carries through a calculation. Every measurement
is wrong to some degree, and a number without an uncertainty is not a
measurement; it is a rumour.

\subsection{The two kinds of error}

Errors come in two kinds that demand \emph{opposite} cures, and telling them apart decides whether
repeating a measurement will help.

\begin{defbox}{Systematic and random errors}
A \textbf{systematic} error biases every reading the same way (a mis-zeroed scale, a bent ruler, an
uncalibrated instrument). A \textbf{random} error scatters readings both ways (a shaky hand, thermal
noise, judgement calls).
\end{defbox}

The crucial fact, constantly missed: \textbf{random errors average away, systematic errors do not.}
Measure a length twenty times with a shaky hand and the mean beats any single reading, because the
scatter is as often positive as negative and cancels. Statistics shows that the uncertainty of the
mean falls as $1/\sqrt N$ for $N$ readings, so a hundred readings pin the mean down ten times better than one
reading does. But if your ruler's zero is worn so every reading runs
\SI{1}{mm} long, then a \emph{million} readings give an exquisitely precise answer that is still
\SI{1}{mm} too long. The scatter shrinks; the bias never moves.

\begin{defbox}{Accuracy and precision: two different words}
\textbf{Accurate} means close to the true value (small \emph{systematic} error).
\textbf{Precise} means that repeated readings are tightly clustered (small \emph{random} error).
They are independent. You can be \emph{precisely wrong}, the most dangerous state in experimental
science, because tight agreement \emph{feels} like success.
\end{defbox}

\why{In 2011 the OPERA experiment reported neutrinos apparently travelling faster than light, with
excellent statistics: the random errors were tiny. A loose fibre-optic connector
was delaying the timing signal by about \SI{73}{ns}, so that every neutrino seemed to arrive early. No amount of data would ever have found it, and none did: someone
eventually checked the cable. Precise, and precisely wrong.}

\subsection{Measuring the uncertainty itself}

Before uncertainties can be combined, each measured quantity needs one. Two sources supply it.

\begin{defbox}{Absolute, relative and percentage error}
Repeat a measurement $n$ times, getting $a_1,a_2,\ldots,a_n$. The best estimate of the true value is the
\textbf{mean} $\bar a=(a_1+\cdots+a_n)/n$. Each reading's \textbf{absolute error} is its departure from
the mean, $\Delta a_i=a_i-\bar a$ (same units as $a$, either sign). Their typical size is the
\textbf{mean absolute error},
\[
\Delta\bar a=\frac{|\Delta a_1|+|\Delta a_2|+\cdots+|\Delta a_n|}{n},
\]
and the result is reported as $\bar a\pm\Delta\bar a$. The \textbf{relative error} $\Delta\bar a/\bar a$ is a
pure number, and multiplied by $100\%$ it is the \textbf{percentage error}; these are what let a length
and a time be compared. A single reading still has an uncertainty: no instrument can resolve less than
its \textbf{least count}, the smallest division it can read (\SI{1}{mm} on a metre rule), so a single
reading is quoted to $\pm$ one least count.
\end{defbox}

Five timings of a pendulum give $1.52$, $1.48$, $1.55$, $1.50$ and \SI{1.45}{s}. The mean is
\SI{1.50}{s}; the absolute errors are $+0.02$, $-0.02$, $+0.05$, $0$ and $-0.05$, whose magnitudes average
to \SI{0.028}{s}. So $T=\SI{1.50 \pm 0.03}{s}$, a relative error of $0.03/1.50=2\%$. Averaging the
magnitudes matters: the signed errors always sum to zero about the mean, which would claim a perfect
measurement.

\subsection{Propagating uncertainty}\label{sec:propagation}

A result computed from several measured quantities inherits their uncertainties. The rules for combining them
are short, and they follow in one line from the logarithmic differentiation of \chref{ch:toolkit}.

\begin{methodbox}{Combining uncertainties}
For a \emph{sum or difference} $Q=x\pm y$, the \emph{absolute} uncertainties add:
$\Delta Q=\Delta x+\Delta y$.
For a \emph{product or quotient} $Q=xy$ or $x/y$, the \emph{relative} uncertainties add:
$\dfrac{\Delta Q}{Q}=\dfrac{\Delta x}{x}+\dfrac{\Delta y}{y}$.
For a power $Q=x^n$, the relative uncertainty multiplies by $|n|$: $\dfrac{\Delta Q}{Q}=|n|\dfrac{\Delta x}{x}$
(the sign of $n$ does not matter: a $1/r^2$ law is as unforgiving as an $r^2$ one).
\end{methodbox}

Even for a difference the absolute errors \emph{add}, since $\Delta x$ and $\Delta y$ may have opposite
signs. This makes the difference of two nearly equal quantities treacherous: with $x=\SI{20.0 \pm 0.1}{cm}$
and $y=\SI{19.0 \pm 0.1}{cm}$, each known to $0.5\%$, $x-y=\SI{1.0 \pm 0.2}{cm}$ is uncertain by $20\%$.

\begin{derivbox}{Where the rule for products comes from}
Take $Q=xy$ and perturb: $Q+\Delta Q=(x+\Delta x)(y+\Delta y)=xy+x\,\Delta y+y\,\Delta x+\Delta x\,\Delta y$.
The last term is a product of two small quantities, negligible. Subtract $Q=xy$ and divide by it:
\[
\frac{\Delta Q}{Q}=\frac{\Delta x}{x}+\frac{\Delta y}{y}.
\]
Equivalently $\ln Q=\ln x+\ln y$, and differentiating gives the rule in one line, which also explains
why powers multiply by $|n|$, since $\ln x^n=n\ln x$. For $Q=x/y$, $\ln Q=\ln x-\ln y$ gives
$dQ/Q=dx/x-dy/y$; since $\Delta y$ may have either sign, the worst case again adds the magnitudes.
\end{derivbox}

\inwords{The syllabus rules \emph{add} the errors, which assumes the worst case: that every error
conspires in the same direction. Truly independent errors combine in \emph{quadrature}
($\Delta Q/Q=\sqrt{(\Delta x/x)^2+(\Delta y/y)^2}$), giving a smaller, more realistic figure: two 3\%
errors give 4.2\%, not 6\%. This is the same quadrature that made the estimation trick work. Use the
adding version for exams, but know it is a deliberately pessimistic bound.}

\begin{trapbox}{powers amplify uncertainty}
Because a power multiplies the relative error by $|n|$, a quantity you cube (like a volume from a
measured length) has \emph{three times} the relative uncertainty of that length. So the quantity that
enters at the highest power is the one to measure carefully.
\end{trapbox}

\begin{exbox}{Which measurement should you improve?}
A sphere's density comes from mass $m=\SI{250 \pm 1}{g}$ and diameter $d=\SI{5.0 \pm 0.1}{cm}$ via
$\rho=m/V$ with $V=\tfrac43\pi(d/2)^3=\pi d^3/6$, so $\rho=6m/\pi d^3$. The relative uncertainties: $\Delta m/m=0.4\%$, but $\Delta d/d=2\%$ enters
\emph{cubed}, contributing $3\times2\%=6\%$. So $\Delta\rho/\rho=0.4\%+6\%=6.4\%$, and
$\rho=\SI{3.8 \pm 0.2}{g/cm^3}$. The lesson: the diameter is \emph{fifteen times} more damaging than the
mass, not because it was measured badly, but because it is cubed. Buying a better balance is pointless;
a calliper twice as precise cuts the total from $6.4\%$ to $3.4\%$. \ans{$\SI{3.8 \pm 0.2}{g/cm^3}$}
\end{exbox}

\section{Summary}\label{sec:summary1}
A dimension records how a quantity is built from mass, length and time, and the principle of
homogeneity says that every term added, subtracted or equated in a physical law must share the same
dimensions. Its two most important uses are these: it can \emph{kill} a wrong formula instantly, and it can
\emph{find} the shape of an unknown law by matching powers of $\mathsf{M}$, $\mathsf{L}$ and $\mathsf{T}$.
What it cannot do is equally important: it never supplies a dimensionless constant, it cannot tell
which terms of a sum are present, it is blind to dimensionless variables such as angles, and it cannot
distinguish quantities that share dimensions (work and torque).

\begin{itemize}
\item $[F]=\mathsf{MLT^{-2}}$, $[E]=\mathsf{ML^2T^{-2}}$, $[G]=\mathsf{M^{-1}L^3T^{-2}}$: derive these from
      defining equations rather than memorising them.
\item Accuracy is closeness to the truth (small systematic error); precision is the tightness of repeated
      readings (small random error). They are independent, and ``precisely wrong'' is the dangerous case.
\item For a sum or difference, \emph{absolute} uncertainties add. For a product or quotient,
      \emph{relative} uncertainties add. A power multiplies the relative uncertainty by $|n|$, so the highest power dominates.
\item Estimate first. An order-of-magnitude answer obtained in ten seconds catches the gross blunders
      (a misplaced power of ten, a wrong unit) made in the following ten minutes.
\end{itemize}

\section{Problem Types}\label{sec:types1}
Measurement questions are among the most predictable in the syllabus, and the six problem types below
cover most of them. Each gives the \emph{tell} that identifies the type from the wording of a
question, the \emph{method} that solves it, and the \emph{trap} that the type is designed to spring.

\archq{Order-of-magnitude estimate}
{``Estimate the number of \ldots'' or ``the order of magnitude of \ldots'', with options a factor of ten
or more apart.}
{Decompose the quantity into factors you can guess to within a factor of two or so, keep one significant
figure throughout, and multiply. Report a power of ten.}
{Over-precision, and slips in the unit conversions (minutes to years) that the estimate is meant to
catch. Write each conversion factor out explicitly.}
{The number of times a human heart beats in a 70-year life is of the order of:}
{$10^{7}$}{$10^{9}$}{$10^{11}$}{$10^{13}$\archans{B}{About $70$ beats per minute, $5\times10^{5}$
minutes per year and $70$ years give $70\times(5\times10^{5})\times70\approx2.5\times10^{9}$, of order
$10^{9}$. Option (A) is what results from counting per hour rather than per minute.}}

\archq{Find the dimensional formula of an unfamiliar constant}
{``Find the dimensions of $a$ (or $b$) in a given equation,'' e.g.\ van der Waals' $\left(P+\frac{a}{V^2}\right)(V-b)=RT$.}
{Use the \textbf{principle of homogeneity}: a quantity added to another must share its dimensions, and
the argument of $\sin$, $\log$ or $e^x$ must be dimensionless. Isolate the unknown and substitute the
dimensions of everything else.}
{Trying to recall the constant. You never need to: the equation itself tells you, because whatever is
added to $P$ must \emph{be} a pressure.}
{In the equation $x=a+bt^2$ ($x$ a distance, $t$ a time), the dimensions of $a/b$ are:}
{$\mathsf{T}$}{$\mathsf{L T^{-2}}$}{dimensionless}{$\mathsf{T^2}$\archans{D}{$a$ is added to $x$, so
$[a]=\mathsf L$. And $bt^2$ must also be a length, so $[b]=\mathsf{LT^{-2}}$. Hence
$[a/b]=\mathsf{L}/\mathsf{LT^{-2}}=\mathsf{T^2}$.}}

\archq{Percentage error in a computed quantity}
{``The percentage errors in $A,B,C$ are given, find the error in $Z=\dfrac{A^pB^q}{C^r}$.''}
{Relative errors \emph{add}, weighted by the powers:
$\dfrac{\Delta Z}{Z}=p\dfrac{\Delta A}{A}+q\dfrac{\Delta B}{B}+r\dfrac{\Delta C}{C}$. The power multiplies
the error whether it is in the numerator or the denominator.}
{Subtracting the denominator's error. Errors accumulate in the \emph{worst case}, so every term carries a
$+$ sign.}
{The percentage errors in $A$, $B$ and $C$ are $1\%$, $2\%$ and $3\%$. The percentage error in
$Z=\dfrac{A^2B^3}{C^4}$ is:}
{$4\%$}{$6\%$}{$20\%$}{$24\%$\archans{C}{$2(1)+3(2)+4(3)=2+6+12=20\%$. Option (A) subtracts the
denominator's term ($2+6-12=-4$) and drops the sign; option (B) ignores the powers. Note $C$ dominates despite its
error being only $3\%$, because of the fourth power.}}

\archq{Error propagated into a derived constant}
{``$g$ is found from a pendulum, or $Y$ from an extension, given the errors in the measured quantities,
find the error in the result.''}
{Rearrange for the wanted quantity \emph{first}, then apply the powers rule. For
$T=2\pi\sqrt{\ell/g}$, invert to $g=\dfrac{4\pi^2\ell}{T^2}$, giving
$\dfrac{\Delta g}{g}=\dfrac{\Delta\ell}{\ell}+2\dfrac{\Delta T}{T}$.}
{Applying the rule to the formula as printed. In $T=2\pi\sqrt{\ell/g}$ the power of $g$ is
$-\tfrac12$, so a careless reading gives half the error instead of twice it.}
{In a pendulum experiment $\ell$ is measured to $2\%$ and $T$ to $3\%$. The percentage error in the
computed $g$ is:}
{$8\%$}{$5\%$}{$11\%$}{$3.5\%$\archans{A}{$g=4\pi^2\ell/T^2$, so
$\dfrac{\Delta g}{g}=\dfrac{\Delta \ell}{\ell}+2\dfrac{\Delta T}{T}=2\%+2(3\%)=8\%$. The period matters
twice as much as the length, which is why it is timed over many oscillations. Option (D),
$2\%+\tfrac12(3\%)$, reads the power of $T$ off the formula as printed, so the period counts half
instead of twice.}}

\archq{Significant figures in a calculation}
{``Give the result of a multiplication/addition to the correct number of significant figures.''}
{Two \emph{different} rules. \textbf{Multiplication/division:} the answer keeps the \emph{fewest
significant figures} among the inputs. \textbf{Addition/subtraction:} the answer keeps the \emph{fewest
decimal places}. Leading zeros never count; trailing zeros after a decimal point do.}
{Applying the multiplication rule to a sum. $12.34+5.6$ is limited by \emph{decimal places} (one), giving
$17.9$, not by significant figures.}
{The product $2.5\times1.24$, to the correct number of significant figures, is:}
{$3$}{$3.1$}{$3.10$}{$3.100$\archans{B}{The arithmetic gives $3.100$, but $2.5$ carries only two
significant figures, so the answer must too: $3.1$.}}

\archq{Checking or deriving a formula by dimensions}
{``Check whether an equation is dimensionally correct,'' or ``deduce how a quantity depends on the
others.''}
{For a \emph{check}, evaluate every term separately and compare. To \emph{derive}, write
$y=C\,a^{p}b^{q}c^{r}$, match the powers of $\mathsf{M}$, $\mathsf{L}$ and $\mathsf{T}$, and solve.
With only three base dimensions in mechanics, at most three unknown exponents can be found, and the
dimensionless constant $C$ never can.}
{Concluding that a dimensionally correct equation is \emph{proved}. The method is blind to pure numbers,
so it can disprove but never fully prove.}
{Which of these is dimensionally \emph{incorrect}? ($v$ speed, $s$ distance, $t$ time, $a$ acceleration)}
{$v^2=u^2+2as$}{$s=ut+\tfrac12at^2$}{$v=u+at^2$}{$s=\tfrac12(u+v)t$\archans{C}{In (C), $at^2$ has
dimensions $\mathsf{LT^{-2}}\cdot\mathsf{T^2}=\mathsf L$, a length, which cannot be added to a
velocity. The correct form is $v=u+at$.}}
